\subsection{问题二分析}

第二问既要判断温度变化，也要比较催化剂组合，还要考察不同组合随温度变化的幅度是否
一致。
这三项任务的比较对象不同：温度影响应在同一组合内部判断，组合差异应在同一温度下
判断，而两者共同变化需要在前两类比较的基础上再作概括。

附件在不同温度下覆盖的组合并不完全相同。若直接比较各温度的全部记录，均值变化会
同时混入温度差异和样本组成变化。为保持比较对象一致，主体分析应使用组合覆盖相同的
温度记录，其他不完整温度只能在具有匹配观测的组合内作补充。同时，题目给出的每个
催化剂组合作为一个完整类别，现有数据不足以把组合差异归因于某个单一配方因素
\cite{searle1980}。

同组配对与同温比较能够回答具体方向和相对高低，但还不能统一衡量组合差异、温度差异
及未被简单相加关系解释的变化。因此，需要在可比数据上建立加和分解。由于每个实验
条件没有重复观测，未解释部分不能单独归因于组合与温度的共同作用或实验误差；模型只作描述性
分解，并通过保持比较结构的删减检查限定结论范围
\cite{nisttwoway,graphpadnorep}。
